\documentclass[10pt,a4paper]{amsart}
\usepackage[margin=19mm]{geometry}
\usepackage{amsmath,amssymb,mathtools}
\usepackage[scr=rsfs,cal=euler]{mathalfa}
\usepackage{xfrac}
\usepackage[unicode,hidelinks,bookmarks=false]{hyperref}
\newcommand{\ie}{i.e.\ }
\newcommand{\cf}{cf.\ }
\newcommand{\resp}{respectively\ }
\newcommand{\Subsection}{\subsection}
\providecommand{\nobreakdash}{\nobreak\hbox{-}}
\setlength{\emergencystretch}{2em}
\multlinegap0pt
\usepackage{amssymb}
\usepackage[shortlabels]{enumitem}
\usepackage{xcolor}
\newtheorem{theorem}{Theorem}
\newcommand{\Mor}{\mathsf{Mor}}
\newcommand{\ODS}{{\mathsf{ODS}}}
\newcommand{\Obj}{\mathsf{Obj}}
\def\cal#1{{\mathcal #1}}
\def\calC{{\cal C}}
\def\calD{{\cal D}}
\def\calF{{\cal F}}
\def\calL{{\cal L}}
\def\calT{{\cal T}}
\def\calX{{\cal X}}
\def\qed{\ifmmode\squareforqed\else{\unskip\nobreak\hfil
\penalty50\hskip1em\null\nobreak\hfil\squareforqed
\parfillskip=0pt\finalhyphendemerits=0\endgraf}\fi}
\def\squareforqed{\hbox{\rlap{$\sqcap$}$\sqcup$}}
\newcommand{\y}{{\mathsf{\bf d}}}
\title{On the completeness of the countable fragment of geometric logic}
\author{Matthew de Brecht}
\date{Source: arXiv:2607.22125v1 (CC BY 4.0)}
\begin{document}
\maketitle

\noindent\emph{Source and licence.} On the completeness of the countable fragment of geometric logic. Version arXiv:2607.22125v1 (CC BY 4.0), distributed by arXiv under the Creative Commons Attribution 4.0 International licence, http://creativecommons.org/licenses/by/4.0/, as recorded in the arXiv licence metadata for this exact version and shown on the abstract page https://arxiv.org/abs/2607.22125v1. Full licence notice: \texttt{licenses/arxiv-2607.22125-CC-BY-4.0.txt}.\par\smallskip

\noindent\textit{Editorial scope.} Counted unit: the article's theorem theorem (source0.tex lines 995--997). The article's complete original proof of it is delivered with it (source0.tex lines 998--1040, one \texttt{proof} environment of 43 lines). No mathematics is written, repaired or re-derived: the statement and the proof are the article's own lines, in the article's own notation, and no retained line is substituted or re-flowed.  Retained uncounted as the source-local context the proof needs: lines 988--991.  Nothing was omitted from any retained span, so the package carries no omission record.  No retained line contains a citation, so the wrapper carries no bibliography and no bibliography entry.  No retained line contains a figure environment, an image-inclusion command or drawing code, so no image is delivered and the package declares no asset dependency.  Editorial formatting only, in addition: an AMS \texttt{amsart} preamble that restates the article's own declarations and macros used by the retained lines, namely the environments \texttt{theorem} on separate counters, together with the standard packages those lines need.  The retained environments print in this document as Theorem 1 in order of appearance, whereas the article prints the same items with the numbering of its own full document.  The complete native source of the article as distributed in the arXiv e-print for this exact version is bundled unchanged as \texttt{sources/205915/source0.tex}.
\begin{theorem}\label{thm:XTsober}
$\y\colon \calX_\calT\to \ODS^{\ODS^{\calX_\calT}}$ restricts to a continuous equivalence between $\calX_\calT$ and the category of $\sigma$-pretopos morphisms from $\ODS^{\calX_\calT}$ to $\ODS$.
\qed
\end{theorem}

\begin{theorem}
Let $\calT_1$ and $\calT_2$ be countable $\sigma$-coherent theories over countable languages $\calL_1$ and $\calL_2$, respectively. There is a dual equivalence between continuous functors $F\colon \calX_{\calT_1} \to \calX_{\calT_2}$ and $\sigma$-coherent morphisms $\calF\colon \ODS^{\calX_{\calT_2}} \to \ODS^{\calX_{\calT_1}}$.
\end{theorem}

\begin{proof}
To simplify notation, set $\calC=\calX_{\calT_1}$ and $\calD=\calX_{\calT_2}$. Given $F\colon \calC\to\calD$ define $\calF_F\colon \ODS^\calD \to \ODS^\calC$ as:
\begin{itemize}
\item
$(\calF_F)_\Obj(G) = G\circ F$, for $G\colon \calD\to\ODS$.
\item
$(\calF_F)_\Mor(\eta\colon G\to H)$ is the natural transformation $\varepsilon \colon G\circ F \to H\circ F$ given by $\varepsilon(C)=\eta(F_\Obj(C))$ for each $C\in\calC_\Obj$.
\end{itemize}
Clearly $\calF_F$ is a $\sigma$-coherent morphism because colimits and limits are computed objectwise.

Next, given $\calF\colon \ODS^\calD \to \ODS^\calC$, first define $\widehat{\calF}\colon \calC\to \ODS^{\ODS^\calD}$ as:
\begin{itemize}
\item
$\widehat{\calF}_\Obj(C)\colon \ODS^{\calD}\to\ODS$ is defined for $C\in\calC_\Obj$ as:
\begin{itemize}
\item
$(\widehat{\calF}_\Obj(C))_\Obj(G) = (\calF_\Obj(G))_\Obj(C)$ for $G\colon \calD\to\ODS$.
\item
$(\widehat{\calF}_\Obj(C))_\Mor(\eta\colon G\to H) = \calF_\Mor(\eta)(C)$.
\end{itemize}
\item
$\widehat{\calF}_\Mor(f\colon C\to C')$ is the natural transformation $\varepsilon\colon \widehat{\calF}_\Obj(C)\to \widehat{\calF}_\Obj(C')$ defined as $\varepsilon(G) = (\calF_\Obj(G))_\Mor(f)$ for each $G\colon \calD\to\ODS$.
\end{itemize}
$\widehat{\calF}_\Obj(C)$ is a $\sigma$-coherent morphism for each $C\in\calC_\Obj$ because $\calF$ is a $\sigma$-coherent morphism and colimits and limits in $\ODS^\calD$ are computed objectwise. Therefore, we can compose $\widehat{\calF}$ with the equivalence from Theorem~\ref{thm:XTsober} to get a continuous functor $F_\calF\colon\calC\to\calD$.

Going from $F\colon \calC\to\calD$ to $\calF_F\colon\ODS^\calD \to \ODS^\calC$ to $\widehat{\calF_F}\colon \calC\to \ODS^{\ODS^\calD}$ as above, we see that $\widehat{\calF_F} = \y\circ F$. Applying the equivalence from Theorem~\ref{thm:XTsober} shows that $F_{\calF_F}$ is naturally isomorphic to $F$.

Conversely, go from $\calF\colon\ODS^\calD \to \ODS^\calC$ to $\widehat{\calF}\colon \calC\to \ODS^{\ODS^\calD}$ to $F_\calF\colon \calC\to\calD$ as above. $\widehat{\calF}$ and $\y\circ F_\calF$ are naturally isomorphic by definition of $F_\calF$. For $G\colon \calD\to\ODS$ and $C\in\calC_\Obj$ and $\eta\colon G\to H$ we have:
\begin{eqnarray*}
(\calF_{F_\calF})_\Obj(G)(C) &=& G_\Obj((F_\calF)_\Obj(C))\\
&=& ((\y\circ F_\calF)_\Obj(C))_\Obj(G)\\
&\cong& (\widehat{\calF}_\Obj(C))_\Obj(G)\\
&=& (\calF_\Obj(G))_\Obj(C)
\end{eqnarray*}
and
\begin{eqnarray*}
(\calF_{F_\calF})_\Mor(\eta)(C) &=& \eta((F_\calF)_\Obj(C))\\
&=& ((\y\circ F_\calF)_\Obj(C))_\Mor(\eta)\\
&\cong& (\widehat{\calF}_\Obj(C))_\Mor(\eta)\\
&=& \calF_\Mor(\eta)(C).
\end{eqnarray*}
Therefore, $\calF_{F_\calF}$ is naturally isomorphic to $\calF$.
\end{proof}

\end{document}
